\question
ООО “Белорусские креветки” базируется в Минске, но для расширения производства приобрело дополнительные помещения в Москве и Санкт-Петербурге. Теперь отдел логистики в ужасе: необходимо что-то решать с цепочками поставок, путь неблизкий, а креветки портятся. Спасать компанию поручили стажёру, то есть - угадайте, кому - вам.

Есть список пунктов, через которые проезжают контейнеры с продукцией. Необходимо, чтобы для каждого пункта число ежедневно прибывающих контейнеров было равно числу увозимых оттуда. Вы собираете информацию по каждому из пунктов на пути в Москву и Петербург, но данные не сходятся, и вам надо найти ошибки в указаниях.

Рассчитайте, какое максимальное количество дней (в идеале) вам понадобится, чтобы доставить партию креветок из Минска в Москву и в Петербург.


Минск-Москва:
\begin{enumerate}
\item Пункт А длительность 1, предшествующие: C,E;
\item Пункт B длительность 2, предшествующие: C,E;
\item Пункт C длительность 1, предшествующие: –;
\item Пункт D длительность 4, предшествующие: B,C;
\item Пункт E длительность 5, предшествующие: С,F;
\item Пункт G длительность 2, предшествующие: A;
\item Пункт F длительность 2, предшествующие: C,G;
\end{enumerate}
Минск-Петербург:
\begin{enumerate}
\item Пункт А длительность 2, предшествующие: E;
\item Пункт B длительность 6, предшествующие: D;
\item Пункт C длительность 2, предшествующие: D,G;
\item Пункт D длительность 4, предшествующие: -;
\item Пункт E длительность 5, предшествующие: F,G;
\item Пункт G длительность 2, предшествующие: -;
\item Пункт F длительность 1, предшествующие: D;
\end{enumerate}
